\section{Purpose of this document}

This document aligns the team on \emph{how} we get to flight / competition:
cold-flow, static fire, packing, on-site days, and countdown. Without it,
assumptions diverge (``when is SF?'', ``what rides in which car?'', ``what on
misfire?'').

\noindent\textbf{Audience:} the whole team --- especially people on the test
stand and at the competition site.

\noindent\textbf{Out of scope:} fastener selection, pinouts, full BOM --- that
belongs in the SAD / subsystem notes.

\subsection{Campaign identification}

\begin{tabular}{@{}ll@{}}
\textbf{Project:} & \ProjectName \\
\textbf{Campaign / mission:} & \MissionName \\
\textbf{Stages in scope:} & \todo{cold-flow, static fire, packing, competition flight} \\
\textbf{Locations:} & \todo{lab / SF stand / competition range} \\
\end{tabular}

\subsection{Campaign goal (one sentence)}

\todo{e.g.\ ``Deliver a recovered competition flight after a successful
cold-flow and static fire.''}

\subsection{Campaign success and failure criteria}

\guidance{Per-stage criteria also live in the cold-flow / SF / flight
sections. Here --- what counts as success for the whole campaign.}

\noindent\begin{tabularx}{\linewidth}{@{}l X@{}}
\toprule
\textbf{Level} & \textbf{Criterion} \\
\midrule
Full success & \todo{e.g.\ competition flight + recovery + data + no injuries} \\
Partial success & \todo{e.g.\ SF passed, flight slipped for weather} \\
Failure / campaign stop & \todo{e.g.\ SF failed, not ready to ship} \\
\bottomrule
\end{tabularx}
